\section{Gestión de estados}
En el Agentic Loop, cada sample puede estar en diferentes estados:
\begin{itemize}
  \item \textbf{Generating}: está generando texto
  \item \textbf{Tool Calling}: está ejecutando una herramienta
  \item \textbf{Waiting}: espera a que la herramienta regrese
  \item \textbf{Done}: ha completado todas las rondas
\end{itemize}

La transición de estados se controla mediante el archivo de configuración: si se configura una tool, se llama a la tool; si se configura una interaction, se simula la retroalimentación del usuario.

\subsection{Resumen de la sección}
La ejecución asíncrona y una gestión de estados fina son fundamentales para el funcionamiento eficiente del Agentic Loop.

\newpage
